\section{Robustness}\label{sec:robustness}

Six dimensions of robustness feed into the headline: the strict-invariance limiting case; the cost-distribution estimator ($F_c^k$ regime); the kernel bandwidth used in the Convite GPV cross-check; the bid-level sample filter; the temporal window; and the set of alternative policy variants considered. The first five test the identification and equilibrium-selection assumptions; the sixth feeds the policy discussion in Section~\ref{sec:discussion}. The central findings survive them: the intensive margin remains dominant in both classes, with shares in the 75--85 percent range, and the welfare comparison remains class-specific rather than universal.

\subsection{Strict-primitive-invariance benchmark}

Table~\ref{tab:v3_strict_invariance} reports the limiting case in which $F_c^{\text{SME,Post}}$ is replaced by $F_c^{\text{SME,Pre}}$ while preserving the observed post-policy SME entry count. This is the cleanest benchmark for a literal primitive-invariance reading of Assumption~\ref{a:setaside}. The decomposition remains intact. The total price effect becomes 0.29 in non-pharma and 0.47 in pharma, with intensive shares of 85 and 79 percent, respectively. Welfare loss at $\lambda = 0.30$ is 30.0 percent of procurement value in non-pharma and 39.2 percent in pharma.

What changes is the policy ranking in the thin, heterogeneous pharma market. In non-pharma, the 10 percent price preference remains Pareto-superior under strict invariance, just as under the main specification. In pharma, the preference ranking reverses and the implied indifference weight falls to $w^{\text{SME}}_\star = 0.7$, versus 3.0 under the main ALS-style equilibrium-selection specification. I read this as mechanism, not fragility: the policy comparison is robust in thick standardized markets and more model-sensitive when entry selection among protected bidders is first-order.

\subsection{Cost-distribution regime}

The baseline all-bidders estimator biases the upper tail of $F_c^k$ slightly downward; losers-only biases it materially upward; a Turnbull NPMLE that treats the winner as left-censored delivers point identification. Table~\ref{tab:v3_sensitivity_fc} runs the BNE decomposition under each. The three point estimates for $p_{S_3} - p_{S_1}$ are:
\begin{center}
\begin{tabular}{lccc}
\toprule
Class      & Losers-only & All-bidders (v3) & Turnbull \\
\midrule
Non-pharma & 0.253 & 0.232 & 0.242 \\
Pharma     & 0.343 & 0.331 & 0.338 \\
\bottomrule
\end{tabular}
\end{center}
All three regimes agree within 5 percent of the all-bidders point. Bootstrap CIs overlap by more than 90 percent across regimes. Turnbull raises the intensive share in pharma from 75 percent (all-bidders) to 81 percent---if anything strengthening the main interpretation. The all-bidders approximation used throughout the paper is tight and defensible without the Turnbull machinery; readers who prefer point identification should see Turnbull and the same decomposition.

\subsection{Kernel bandwidth}

The cross-modality check in Section~\ref{sec:identification} uses GPV inversion on Convite bids, which requires a kernel density estimate $g(b)$. Silverman's rule-of-thumb $h = 1.06 \hat\sigma n^{-1/5}$ is the baseline. Table~\ref{tab:v3_bandwidth_grid} runs the GPV under $h$-factors of 0.5, 0.75, 1.0, 1.5, and 2.0. The resulting range of $c_{0.5}$ (median cost) across the eight (period $\times$ pharma $\times$ type) strata is 1.2--3.5 percent of the baseline in seven strata and 5.2 percent in the thinnest cell (pharma SME Pre, $n = 2{,}300$). This is second-order to the BNE Monte Carlo noise of 1--2 percent. The $c_{0.75}$ range is 0.1--1.2 percent across all eight strata. Bandwidth is not a first-order concern for the identification.

\subsection{Bid-level sample filter}

The baseline filter is $c_\epsilon \le 3$ and $n \ge 2$. Table~\ref{tab:v3_filter_sensitivity} compares four alternatives:
\begin{itemize}
\item \emph{Tight} ($c_\epsilon \le 2$, $n \ge 3$): $p_{S_3} - p_{S_1}$ is 7 percent lower in non-pharma (0.213 vs.\ 0.228) and 11 percent lower in pharma (0.278 vs.\ 0.312).
\item \emph{Very tight} ($c_\epsilon \le 1.5$, $n \ge 3$): 15 and 19 percent lower respectively. The intensive share in pharma jumps to 99 percent---tightening the cost distribution to its central bulk effectively removes the tail where the entry response operates.
\item \emph{Loose} ($c_\epsilon \le 5$): 3 percent higher in non-pharma and 27 percent higher in pharma. The pharma move is concentrated in the extra 1{,}000 observations with $c_\epsilon \in (3, 5]$---almost always outliers at the 500--1000 percent-of-reference band, probably data-entry errors at the hundred-times level. The baseline cut is designed to truncate these.
\end{itemize}
The baseline sits between tight and loose in both classes, is inside the bootstrap CI under any reasonable cut, and can be defended directly: the 3x-reference upper bound catches plausible error patterns while preserving 99.6 percent of the data.

\subsection{Temporal window}

Table~\ref{tab:v3_window_sensitivity} varies the symmetric window around the March 2018 cutoff: 18 months (baseline), 12 months, and 6 months. The total effect $p_{S_3} - p_{S_1}$ is 0.218 (non-pharma 18m), 0.250 (12m), and 0.282 (6m); 0.316 (pharma 18m), 0.295 (12m), 0.374 (6m). Non-pharma is monotonic in window length; pharma has a non-monotonic dip at 12m that is likely a boundary-effect artifact (CMED regulatory events in Q1 and Q4 2018 that partially contaminate the 12m window without full inclusion in the 18m extent).

The more notable finding is on the intensive share:
\begin{center}
\begin{tabular}{lccc}
\toprule
Class      & 18m    & 12m    & 6m \\
\midrule
Non-pharma & 68.1\% & 76.4\% & 77.1\% \\
Pharma     & 67.0\% & 69.4\% & 78.1\% \\
\bottomrule
\end{tabular}
\end{center}
The intensive share rises as the window shrinks. The reading is that the entry response is gradual---SMEs take time to observe the regime, register as eligible suppliers, and start bidding---so narrower windows capture the immediate bidding response before the extensive-margin adjustment has equilibrated. The 18m specification reports the mature, steady-state number; shorter windows would oversell the intensive margin. This is a substantive finding and arguably strengthens the case for the 18m baseline as the appropriate horizon.

\subsection{Alternative policy variants}

The 10 percent price preference in Section~\ref{sec:pareto} is informative partly because the alternative variants reveal what bidder exclusion is buying. Table~\ref{tab:v3_apv} also reports two intermediate variants that serve as policy-design sensitivities:

\begin{itemize}
\item \emph{V1} (partial set-aside): 50 percent of auctions run SME-only at the Post pool, 50 percent open at the Pre pool. $\Delta p$ is 30--45 percent of V0's---roughly half, but not exactly, reflecting the non-linearity of pool restriction. A partial rule buys roughly proportional fiscal savings but at the cost of a more complex administrative rule.
\item \emph{V2} (entry-isolated counterfactual): full set-aside with the SME pool held at pre-period size (no entry response). $\Delta p$ is 162--170 percent of V0. Without endogenous entry the policy would be materially more costly; the extensive-margin response absorbs 60--70 percent of what would otherwise be a heavier intensive shock. This is why Section~\ref{sec:results} treats entry as the government's partial insurance.
\end{itemize}

Taken together, V1--V3 span the policy design space between the status-quo full set-aside and the open regime. V2 isolates what the entry response is worth; V1 prices a middle-ground design; V3 shows that a well-targeted preference instrument captures the SME-protection value at essentially zero price cost in thick markets and gives a diagnostic benchmark in thin markets.
